\section{おわりに}

本稿では，TPL検出器が生成した関数候補を集約前に再検証する手法を示した．GhidraのHigh P-codeから既定契約上の戻り値と効果を構成し，SMTソルバで差を調べる．観測対象の経路，CALL，メモリ，入出力，終了またはループを十分に表現できない場合は，差が見つからなくてもEQとせずUNKNOWNとする．

EqBenchの4関数・480ペアでは，257ペアがEQ，216ペアがUNKNOWN，7ペアがTIMEOUTであった．最適化，コンパイラ，ISAを個別に層別したうえで，4関数全体をループ・CALLの有無でも集計した．解析版A内では，\texttt{ltfive/lib}のUNKNOWNに戻り値観測契約・幅の由来不足が，\texttt{ltfive/client}のUNKNOWNにCALL位置対応と経路被覆・反例安全性が多く現れた．解析版Bの\texttt{LoopMult10/foo}では，非循環45ペアがすべてEQで，循環を含む75ペアがすべてUNKNOWNとなり，ループ要約と反例安全性の未完了が対応した．これらはペア内で重複する原因の該当数であり，版AとBをまたぐ効果量の比較ではない．また，7件のTIMEOUTは\texttt{ltfive/client}に集中したが，保存ログから処理段階やCALLとの因果は特定できない．非等価ペアと実際のTPL候補への適用は今後の課題である．
